\documentclass[../../main2.tex]{subfiles}

\begin{document}

	% ======================================================================
	% 骨架文件：小节标题与追问链为终版；正文由后续会话填充。
	% 硬要求（notes/arena-prompt.md 任务卡 D）：
	%   §13.3 必须真实推导一个具体例子（不许空话）；
	%   §13.4 失败模式清单必须列满 5 条以上。
	% ======================================================================

	\chapter{从贡献表到未来走向}
	\label{ch:forecast}

	\begin{motivation}
		上一章遗留的问题：四条通道全部定义完毕，贡献分解的公式在手——但那是对\textbf{过去}的分解。我们手里拿着 $C_{\text{模因}}, C_{\text{经济}}, C_{\text{政治}}, C_{\text{法律}}$ 四个数字，怎么用它说出「接下来会怎样」？\\
		\textbf{核心动机}：把贡献表拼回时间线——重组条件概率链，给出一条可辩护的未来。这是全书唯一一次完整演示「如何用这本书」。
	\end{motivation}

	\begin{logicmap}
	\begin{tikzpicture}[node distance=0.9cm, every node/.style={draw, rectangle, rounded corners, align=center}]
	\node (q) {贡献表怎么变成「接下来」？};
	\node (s) [below=of q] {贡献表 = 状态压缩（§13.1）};
	\node (a) [below=of s] {条件概率链重组（§13.2）};
	\node (e) [below=of a] {一个完整演示（§13.3）};
	\node (f) [below=of e] {失败模式清单（§13.4）};
	\draw[-{Stealth}] (q) -- (s);
	\draw[-{Stealth}] (s) -- (a);
	\draw[-{Stealth}] (a) -- (e);
	\draw[-{Stealth}] (e) -- (f);
	\end{tikzpicture}
	\end{logicmap}

	% ==================================================================
	\section{贡献表是一份状态压缩}
	\label{sec:fc-table}
	把四个通道的贡献 $C_{\text{模因}}, C_{\text{经济}}, C_{\text{政治}}, C_{\text{法律}}$ 放在一起，就得到一张状态表。这张表是对当前社会状态的粗粒度压缩：它告诉你「此刻各通道在多大程度上驱动变化」，但不告诉你细节。

	丢了什么？丢个体差异（一亿人的行为被压成四个数字）、丢通道内部异质性（经济通道里，房地产和互联网对同一个指标的贡献方向可能相反，但表里只有一个数）。这和第2章的地图类比同构：地图不是领土，但好的地图足以导航。能接受多少丢失，取决于预测目标的时间尺度与粒度——预测明年 GDP 增速，四个数字够了；预测某个小区房价，不够。

	\begin{readingnote}
		天气预报不说「每一滴雨落在哪里」。它说「明天降水概率 70%，降水量 10-20mm」。这是同一套压缩逻辑：用少量状态变量描述一个高维系统，牺牲分辨率换取可预测性。我们的四通道贡献表就是社会系统的「降水概率」。
	\end{readingnote}

	\paragraph*{逻辑衔接}
	压缩后的状态要能重新展开成未来。下一节做这个展开。
	% ==================================================================
	\section{条件概率链的重新组装}
	第6章建立了乘积分解：$P(F) = \prod_k P(C_k \mid C_{k-1})$——未来的概率等于各条件因子概率的连乘。这里把抽象的因子换成具体的通道贡献：每个 $P(C_k \mid C_{k-1})$ 对应「在第 k-1 步状态已知时，第 k 通道贡献落在某区间的概率」。

	和第6章的区别在于可执行性。第6章是概念框架——告诉你「应该这么分解」。这一章给出测量接口：每个因子的数据来源是什么（统计年鉴、调查数据、政策文本）、更新规则是什么（新数据出来怎么修正概率）、误差怎么传播（一个因子估错了，最终结果偏多少）。附录 A.2 给出了完整的数学接口定义。

	流程：(1) 从贡献表读取当前四通道状态 → (2) 对每个通道估计下一步的条件概率分布 → (3) 连乘得到联合分布 → (4) 从联合分布提取断言（最可能值、置信区间、尾部风险）。
	% ==================================================================
	\section{一个具体工作过的例子}
	\label{sec:fc-example}
	选一个读者熟悉的事件：某城市宣布取消限购政策后，半年内房价上涨的概率是多少？

	第一步：状态刻画。当前该城市住房市场处于供大于求状态（库存去化周期 18 个月），居民杠杆率接近上限，政策预期已经部分 price-in（消息公布前一周成交量已上升 30%）。

	第二步：四通道贡献估计。$C_{\text{模因}}$：政策信号通过媒体和社交网络扩散，购房信心提升——估计贡献 +15-25%（基于类似政策的历史传播速度）。$C_{\text{经济}}$：供需基本面未变（库存仍高），但流动性边际改善——估计贡献 +5-10%。$C_{\text{政治}}$：地方政府有维稳动机，可能配合出台补贴——估计贡献 +5-15%（不确定性大）。$C_{\text{法律}}$：无直接法律约束变化——贡献 ≈ 0。

	第三步：链式组装。$P(\text{涨价} \mid \text{当前状态}) \approx P(C_{\text{模因}} > \theta_1) \times P(C_{\text{经济}} > \theta_2) \times P(C_{\text{政治}} > \theta_3)$，其中 $\theta_i$ 是各通道触发涨价的阈值。代入估计值：$0.6 \times 0.4 \times 0.5 = 0.12$。

	第四步：概率断言。半年内房价上涨的概率约 12%，95% 置信区间 [5%, 25%]。依赖因素清单：如果 $C_{\text{政治}}$ 的贡献被低估（实际出台了强力补贴），概率上修至 ~30%；如果全国经济下行拖累需求，$C_{\text{经济}}$ 转负，概率下修至 ~5%。\lowfreq{例子的具体数字是示意性估计，不是实证结果。}

	\begin{questionbox}
		「一个例子证明不了方法有效。」对。例子演示的是可执行性——你能不能真的算出一个数字。有效性要看失败模式清单（下一节）与方法边界。一个跑通的例子只能说明「没有卡住」，不能说明「算对了」。
	\end{questionbox}

	\paragraph*{逻辑衔接}
	演示跑通了。现在做另一半诚实的事：这个方法会在哪里翻车？
	% ==================================================================
	\section{预测失败模式清单}
	\label{sec:fc-failure}
	五种失败模式，按严重程度排序：

	(1) **权重不可测。** 症状：某个通道的贡献无法从现有数据中分离出来（例如 $C_{\text{政治}}$ 在威权体制下几乎不产生可观测信号）。为什么无法修复：不是数据不够多，而是因果通道本身被制度屏蔽了。部分缓解：用代理变量（如官员更替频率）间接估计，但精度大幅下降。

	(2) **交互项过大。** 症状：两个通道的联合效应远大于各自单独效应之和（例如经济危机 × 政治极化 → 社会动荡，单独看都不够触发）。为什么无法修复：乘积分解假设通道近似独立，强耦合时这个假设失效。部分缓解：加入二阶交互项，但参数数量爆炸，需要更多数据。

	(3) **通道漏项。** 症状：出现了第四通道之外的驱动力量（例如一场瘟疫同时改变了经济、政治、文化三个通道，但它不属于任何一条）。为什么无法修复：通道分类是基于当前知识划分的，新现象可能超出分类框架。部分缓解：保留「残差通道」吸收未解释方差，但残差越大说明模型越弱。

	(4) **外部冲击。** 症状：系统受到来自分类框架之外的冲击（战争、小行星撞击、技术奇点）。为什么无法修复：冲击的定义就是「不可预测」，任何预测框架对此都无能为力。部分缓解：压力测试——模拟极端情景下的系统响应，但不改变基准预测。

	(5) **非平稳。** 症状：通道的动力学参数随时间漂移（例如模因传播速度因互联网出现而突变）。为什么无法修复：历史数据训练出的参数在新环境下失效。部分缓解：滑动窗口更新 + 参数漂移监测，但响应永远滞后于变化。

	少于 3 条的失败清单不值得信任。一个方法如果只列出「数据不够」和「模型简化」两条失败原因，那它要么没认真想，要么在回避真正的弱点。

	\begin{reader_anticipation}
		\item 失败之后怎么办？→ 第14章：预测的反身性——预测本身改变被预测的世界。
	\end{reader_anticipation}

	\paragraph*{逻辑衔接}
	我们会算未来了，也知道了哪里算不动。还剩最后一问：我们预测出的未来，会不会反过来改变人的动机？
	% ==================================================================
	\section{本章小结与衔接}
	\label{sec:fc-summary}
	从贡献表到概率断言的完整链路：状态压缩（四通道贡献表）→ 条件概率链组装（乘积分解 + 测量接口）→ 概率断言（最可能值 + 置信区间 + 依赖因素清单）。边界明确：五种失败模式（权重不可测、交互项过大、通道漏项、外部冲击、非平稳）定义了方法的适用域。

	\paragraph*{逻辑衔接}
	第14章「循环的闭合：未来塑造新的动机」——全书最后一章，回到第1章的起点。
\end{document}
